\documentclass[11pt]{article}
\usepackage[a4paper,margin=1in]{geometry}
\usepackage{amsmath,amssymb,amsthm}
\usepackage{graphicx}
\usepackage{xcolor}
\usepackage[hidelinks]{hyperref}
\usepackage{microtype}
\setlength{\emergencystretch}{2em}

\newtheorem{theorem}{Theorem}
\newtheorem{corollary}[theorem]{Corollary}
\newtheorem{remark}[theorem]{Remark}
\newcommand{\R}{\mathbb R}
\newcommand{\one}{\mathbf 1}
\newcommand{\cM}{\mathcal M}
\newcommand{\Cl}{\mathrm{Cl}}

\title{A missing sharpness edge for spherical maximal functions\\
on H-type groups with center dimension three or seven}
\author{Automated mathematical discovery run \texttt{fa\_banach\_001}}
\date{13 August 2026}

\begin{document}
\maketitle

\begin{center}
\fcolorbox{black}{gray!8}{\parbox{0.91\linewidth}{
\textbf{Status.} Candidate substantial partial result, likely valid, subject
to human review.  We prove the missing $Q_3Q_4$ necessary inequality for
untilted H-type groups whose center has dimension $m=3$ or $m=7$, and more
generally under an explicit common-invariant-block hypothesis.  General
M\'etivier groups and arbitrary tilts are not resolved.
}}
\end{center}

\section{Source problem}

The source is J.~Roos, A.~Seeger, and R.~Srivastava,
\emph{Lebesgue space estimates for spherical maximal functions on Heisenberg
groups}, arXiv:2103.09734v2~\cite{RSS}.  In the notation of that paper, the
horizontal dimension is $h=2n$, the center dimension is $m$, and the total
Euclidean dimension is
\[
 D=h+m.
\]
For skew-symmetric $h\times h$ matrices $J_1,\ldots,J_m$, write
$x^\intercal\vec J\omega=(x^\intercal J_1\omega,\ldots,
x^\intercal J_m\omega)$.  If $\Lambda:\R^h\to\R^m$ is the tilt, the local
maximal operator considered in the source is
\begin{equation}\label{eq:maximal}
 \cM f(x,z)=\sup_{1\leq t\leq2}\left|
 \int_{S^{h-1}}f\bigl(x-t\omega,
 z-tx^\intercal\vec J\omega-t^2\Lambda\omega\bigr)\,d\sigma(\omega)
 \right|.
\end{equation}

The source proves positive estimates in a quadrilateral with vertices
\begin{align}\label{eq:q3q4}
 Q_3&=\left(\frac{D-1}{D+m},\frac{m+1}{D+m}\right),\\
 Q_4&=\left(\frac{D(D-1)}{\Delta},
 \frac{(m+1)(D-1)}{\Delta}\right),
 \qquad \Delta=D^2+(D+1)m+1.\nonumber
\end{align}
Its Remark~1.2(iv), on source p.~5, asks for analogues when $m\geq2$ of its
new example proving sharpness of the line $Q_3Q_4$ for $m=1$:

\begin{center}
\includegraphics[width=\linewidth]{figures/open_problem_crop.png}
\end{center}

\noindent In particular, it asks whether similar examples exist for
$m\geq2$ and whether the corresponding sharpness problem can be settled.

\section{Proof intuition}

The source's $m=1$ construction singles out a two-plane invariant under the
single complex structure $J$.  Horizontal directions in that plane are
localized at scale $\delta$, directions in its orthogonal complement at scale
$\sqrt\delta$, and the center at scale $\delta$.  The apparent
$\sqrt\delta$ error in the group law becomes a product of two
$\sqrt\delta$ errors, hence has size $\delta$.

The same mechanism works for center dimension $m$ if all $J_i$ preserve a
common $(m+1)$-dimensional horizontal block.  The resulting volume exponents
give exactly the line through the published $Q_3$ and $Q_4$.  H-type groups
are real Clifford modules.  In the exceptional center dimensions $3$ and
$7$, their irreducible Clifford blocks have dimensions $4$ and $8$---exactly
$m+1$---so the construction applies.  The block hypothesis is also the
limitation: a general skew pencil need not have such a subspace, and a general
tilt can see the coarse $\sqrt\delta$ directions.

\section{Invariant-block Knapp theorem}

\begin{theorem}\label{thm:block}
Suppose there is a common invariant subspace $V\subset\R^h$ for
$J_1,\ldots,J_m$ with dimension $r$.  Assume also
\begin{equation}\label{eq:tiltcondition}
 \Lambda|_{V^\perp}=0.
\end{equation}
If \eqref{eq:maximal} is bounded from $L^p(\R^{h+m})$ to
$L^q(\R^{h+m})$, then
\begin{equation}\label{eq:generalnec}
 D-m+r-2+\frac{D+m-r}{q}\ \geq\
 \frac{D+m+r}{p}.
\end{equation}
In particular, if $r=m+1$, boundedness forces
\begin{equation}\label{eq:q3q4line}
 (D-1)\left(1+\frac1q\right)\geq
 \frac{D+2m+1}{p}.
\end{equation}
Equality in \eqref{eq:q3q4line} is precisely the line through the two points
$Q_3,Q_4$ in \eqref{eq:q3q4}.
\end{theorem}

\begin{proof}
Let $P$ and $P^\perp$ be the orthogonal projections onto $V$ and
$V^\perp$.  Since every $J_i$ is skew-symmetric and preserves $V$, it also
preserves $V^\perp$.  Fix a small $\delta>0$ and, with a sufficiently large
constant $C$, set
\[
 E_\delta=\{(y,w): |Py|\leq C\delta,\quad
 |P^\perp y|\leq C\delta^{1/2},\quad |w|\leq C\delta\},
 \qquad f_\delta=\one_{E_\delta}.
\]
Then
\begin{equation}\label{eq:inputnorm}
 \|f_\delta\|_p\asymp
 \delta^{\{(h-r)/2+r+m\}/p}
 =\delta^{(D+m+r)/(2p)}.
\end{equation}

Choose a relatively compact coordinate patch $U$ of the unit sphere in $V$.
Consider the output set
\begin{multline}\label{eq:outputset}
 R_\delta=\{(x,z): 9/8\leq |Px|\leq15/8,\quad
 Px/|Px|\in U,\\
 |P^\perp x|\leq c\delta^{1/2},\quad
 |z-|Px|\Lambda x|\leq c\delta\},
\end{multline}
where $c>0$ is a sufficiently small fixed constant.  Fubini's theorem in
orthogonal coordinates gives
\begin{equation}\label{eq:outputvolume}
 |R_\delta|\asymp \delta^{(h-r)/2+m}
 =\delta^{(D+m-r)/2}.
\end{equation}

Fix $(x,z)\in R_\delta$, put $t=|Px|$, and let $v=Px/t$.  In local sphere
coordinates define $\Omega_x\subset S^{h-1}$ by requiring
\[
 |P^\perp\omega|\leq c\delta^{1/2}
 \quad\text{and}\quad
 \operatorname{dist}_{S(V)}
 \left(\frac{P\omega}{|P\omega|},v\right)\leq c\delta.
\]
There are $h-r$ independent coordinates of size $\delta^{1/2}$ and $r-1$
independent coordinates of size $\delta$.  Hence
\begin{equation}\label{eq:capmeasure}
 \sigma(\Omega_x)\asymp
 \delta^{(h-r)/2+r-1}
 =\delta^{(D-m+r-2)/2}.
\end{equation}
Moreover, for $\omega\in\Omega_x$, the unit-sphere equation gives
\begin{equation}\label{eq:horizontalerrors}
 |P(x-t\omega)|\lesssim\delta,
 \qquad |P^\perp(x-t\omega)|\lesssim\delta^{1/2}.
\end{equation}

For each $i$, skew-symmetry yields
\[
 x^\intercal J_i\omega=(x-t\omega)^\intercal J_i\omega.
\]
Because $V$ and $V^\perp$ are invariant, the mixed terms in the orthogonal
decomposition vanish.  Thus \eqref{eq:horizontalerrors} gives
\begin{align*}
 |x^\intercal J_i\omega|
 &\leq |P(x-t\omega)|\,\|J_i\|\,|P\omega|\\
 &\quad+|P^\perp(x-t\omega)|\,\|J_i\|\,|P^\perp\omega|
 \lesssim\delta.
\end{align*}
Using $t=|Px|$ and \eqref{eq:tiltcondition}, the vertical argument in
\eqref{eq:maximal} can be rearranged as
\begin{align*}
 z-tx^\intercal\vec J\omega-t^2\Lambda\omega
 &=z-t\Lambda x+t\Lambda(x-t\omega)
   -t x^\intercal\vec J\omega,
\end{align*}
and every term on the right is $O(\delta)$.  Enlarging $C$ if necessary,
the full argument of $f_\delta$ lies in $E_\delta$ for all
$\omega\in\Omega_x$.  Consequently
\begin{equation}\label{eq:lowerbound}
 \cM f_\delta(x,z)\gtrsim\sigma(\Omega_x)
 \asymp\delta^{(D-m+r-2)/2},
 \qquad (x,z)\in R_\delta.
\end{equation}

Combining \eqref{eq:inputnorm}, \eqref{eq:outputvolume}, and
\eqref{eq:lowerbound}, an $L^p\to L^q$ estimate independent of $\delta$
requires
\[
 \frac{D-m+r-2}{2}+\frac{D+m-r}{2q}
 \geq\frac{D+m+r}{2p},
\]
which is \eqref{eq:generalnec}.

If $r=m+1$, this becomes \eqref{eq:q3q4line}.  Direct substitution of
\eqref{eq:q3q4} gives equality at both $Q_3$ and $Q_4$; equivalently, the
equality line is
\[
 (D-1)(1+1/q)=(D+2m+1)/p.
\]
\end{proof}

\section{H-type consequence}

Recall that for an H-type group one may choose an orthonormal basis of the
center so that
\begin{equation}\label{eq:clifford}
 J_iJ_j+J_jJ_i=-2\delta_{ij}I.
\end{equation}
Thus the horizontal space is a finite-dimensional real module over the
Clifford algebra $\Cl_{0,m}$; see Kaplan~\cite{Kaplan}.

\begin{corollary}\label{cor:htype}
Let $G$ be an H-type group with center dimension $m=3$ or $m=7$.
For the untilted spherical maximal operator ($\Lambda=0$), $L^p\to L^q$
boundedness implies the exact $Q_3Q_4$ necessary inequality
\eqref{eq:q3q4line}.

The same conclusion holds for a tilt $\Lambda$ whenever there is an
irreducible Clifford block $V$ such that
$\operatorname{range}\Lambda^\intercal\subset V$.
\end{corollary}

\begin{proof}
The real Clifford-algebra classification gives
\[
 \Cl_{0,3}\cong\mathbb H\oplus\mathbb H,
 \qquad
 \Cl_{0,7}\cong M_8(\R)\oplus M_8(\R).
\]
Accordingly, every finite-dimensional real module is an orthogonal direct sum
of invariant irreducible modules of real dimension $4$ when $m=3$ and $8$
when $m=7$~\cite{LawsonMichelsohn}.  Choose one such block $V$.  In both
cases $\dim V=m+1$.  For $\Lambda=0$, condition
\eqref{eq:tiltcondition} is automatic.  More generally,
$\operatorname{range}\Lambda^\intercal\subset V$ is equivalent to that
condition.  Theorem~\ref{thm:block} applies.
\end{proof}

Together with the other necessary conditions already proved in the source,
Corollary~\ref{cor:htype} settles the previously missing geometric sharpness
edge for these H-type subcases.  It does not decide the two strong-type
endpoint questions already left open by the source's positive theorem.

\section{Verification, limitations, and novelty status}

The reusable script \texttt{code/verify\_exponents\_and\_clifford.py}
symbolically checks the $Q_3,Q_4$ substitutions and exactly checks concrete
quaternionic and octonionic signed-permutation matrices satisfying
\eqref{eq:clifford}.  These finite checks support, but do not replace, the
proof.  The accompanying \texttt{verification.md} records the command and
scope audit.

The construction requires an invariant $(m+1)$-block and a tilt annihilating
its orthogonal complement.  For a general M\'etivier pencil,
$\operatorname{span}\{u,J_1u,\ldots,J_mu\}$ need not be invariant because
the vectors $J_iJ_ju$ introduce new directions.  A moving-block attempt meets
the same obstruction as a failure of involutivity of the distribution
$\operatorname{span}\{J_1u,\ldots,J_mu\}$ on the sphere.  No claim beyond
the stated hypotheses is made.

A bounded search of the run indexes and source corpus used the exact source
id and phrases involving $Q_3Q_4$, H-type groups, M\'etivier sharpness, and
center dimension $m\geq2$.  The later general two-step result
arXiv:2309.07725~\cite{RyuSeeger} was also checked; it treats the diagonal
global maximal problem and still describes the off-diagonal improving
problem for $m\geq2$ as complicated.  Two external exact-phrase searches
found no statement of the theorem above.  This is a bounded check only, so
novelty confidence remains provisional.

\medskip
\noindent\textbf{Human-review recommendation.}
Check first the cap exponent \eqref{eq:capmeasure}, the vanishing of the mixed
bilinear terms, and the cited real Clifford-module dimensions.  If those pass,
the result appears ready as a substantial partial answer to Remark~1.2(iv).

\begin{thebibliography}{9}

\bibitem{RSS}
J.~Roos, A.~Seeger, and R.~Srivastava,
\emph{Lebesgue space estimates for spherical maximal functions on Heisenberg
groups}, Int. Math. Res. Not. IMRN 2022, no.~24, 19222--19257;
arXiv:2103.09734v2.

\bibitem{Kaplan}
A.~Kaplan,
\emph{Fundamental solutions for a class of hypoelliptic PDE generated by
composition of quadratic forms}, Trans. Amer. Math. Soc. 258 (1980), 147--153.

\bibitem{LawsonMichelsohn}
H.~B.~Lawson, Jr. and M.-L.~Michelsohn,
\emph{Spin Geometry}, Princeton Mathematical Series 38, Princeton University
Press, 1989, Chapter I.

\bibitem{RyuSeeger}
J.~Ryu and A.~Seeger,
\emph{Spherical maximal functions on two step nilpotent Lie groups},
arXiv:2309.07725.

\end{thebibliography}

\end{document}
